\subsection{Nuages de mots}

\begin{figure}[H]
\centering
\includegraphics[width=\textwidth]{figures/06_wordclouds.png}
\caption{Nuages de mots après prétraitement : tweets négatifs (gauche, rouge) et positifs (droite, vert).}
\label{fig:wc}
\end{figure}

Les mots dominants confirment la cohérence des annotations :
\textit{sad, hate, miss, tired, cry} pour le négatif ;
\textit{happy, love, thank, great, amazing} pour le positif.

\subsection{Courbes d'apprentissage LSTM}

\begin{figure}[H]
\centering
\includegraphics[width=0.95\textwidth]{figures/06_lstm_curves.png}
\caption{Évolution de la perte et de la précision du LSTM Bidirectionnel (train vs validation).}
\label{fig:lstm_curves}
\end{figure}

Le modèle converge en 5--8 époques. L'écart faible entre train et validation confirme
l'absence de sur-apprentissage grâce aux couches Dropout.

\subsection{Courbes d'apprentissage BERT}

\begin{figure}[H]
\centering
\includegraphics[width=0.95\textwidth]{figures/06_bert_curves.png}
\caption{Évolution de la perte et de la précision pendant le fine-tuning de DistilBERT (3 époques).}
\label{fig:bert_curves}
\end{figure}

BERT converge très rapidement en 3 époques seulement, grâce aux représentations
pré-entraînées sur des milliards de mots.

\subsection{Projection t-SNE des embeddings Word2Vec}

\begin{figure}[H]
\centering
\includegraphics[width=0.9\textwidth]{figures/06_tsne_w2v.png}
\caption{Projection t-SNE en 2D de 300 mots fréquents (Word2Vec, 200 dimensions).
Les mots de sentiment similaire forment des clusters distincts.}
\label{fig:tsne}
\end{figure}

La projection t-SNE révèle que les mots positifs (\textcolor{posgreen}{vert}) et négatifs
(\textcolor{negred}{rouge}) se regroupent dans des zones distinctes de l'espace sémantique,
confirmant la qualité des embeddings Word2Vec appris.

\subsection{Comparaison finale}

\begin{figure}[H]
\centering
\includegraphics[width=0.9\textwidth]{figures/06_final_comparison.png}
\caption{F1-score final de tous les modèles avec la représentation vectorielle utilisée.}
\label{fig:final}
\end{figure}
